\documentclass[11pt,a4paper]{article}
\usepackage[utf8]{inputenc}
\usepackage[T1]{fontenc}
\usepackage[french]{babel}
\frenchbsetup{StandardLists=true}
\usepackage[left=2cm,right=2cm,top=2cm,bottom=2cm]{geometry}
\usepackage[dvipsnames]{xcolor}
\usepackage{fouriernc}
\usepackage{amsmath,amssymb,amsfonts}
\usepackage{array,tabularx,booktabs,multirow}
\usepackage{colortbl}
\usepackage{graphicx}
\usepackage{mdframed}
\usepackage{enumitem}
\usepackage{fancyhdr}
\usepackage{chemformula}
\usepackage{awesomebox}
\setlength{\aweboxvskip}{0mm}
\usepackage{tcolorbox}
\tcbuselibrary{skins}
\usepackage{fontawesome5}

\definecolor{exoheader}{HTML}{1E5AA0}
\definecolor{exolight}{HTML}{DCE8FF}
\definecolor{warnorange}{RGB}{220,100,0}
\definecolor{problemebg}{HTML}{FFFDE7}

\renewcommand{\arraystretch}{1.8}
\setlength{\extrarowheight}{1mm}

\pagestyle{fancy}
\renewcommand\headrulewidth{1pt}
\fancyhead[L]{Académie de Besançon}
\fancyhead[C]{\textbf{TP --- Titrage de l'ibuprofène}}
\fancyhead[R]{Terminale / BTS Chimie}
\fancyfoot[C]{page \thepage}
\fancyfoot[R]{2025 -- 2026}
\fancyfoot[L]{Alexis RUIZ}

\newcommand{\sectitle}[1]{%
  \vspace{6pt}
  \noindent\colorbox{exoheader}{\parbox{\dimexpr\linewidth-2\fboxsep}%
    {\color{white}\bfseries\large\strut\quad #1}}
  \vspace{4pt}
}
\newcommand{\subsectitle}[1]{%
  \vspace{4pt}
  \noindent\colorbox{exolight}{\parbox{\dimexpr\linewidth-2\fboxsep}%
    {\color{exoheader}\bfseries\strut\quad #1}}
  \vspace{2pt}
}
\newcounter{question}
\newcommand{\question}{%
  \stepcounter{question}%
  \noindent\textcolor{exoheader}{\textbf{Question \thequestion.}}\quad
}

\begin{document}

\noindent\textbf{NOM :}\hfill\textbf{Prénom :}\hfill\textbf{Classe :}\hfill\phantom{x}\\[0.2cm]
\hrule\vspace{0.2cm}

\begin{center}
  {\large\bfseries\textcolor{exoheader}{Travaux Pratiques de Chimie}}\\[4pt]
  {\LARGE\bfseries Titrage acido-basique de l'ibuprofène}\\[2pt]
  {\small Académie de Besançon --- Belfort}
\end{center}
\hrule\vspace{6pt}

\begin{tcolorbox}[enhanced, colback=problemebg, colframe=exoheader, boxrule=1.5pt,
  title={\faFlask\quad Problématique},
  fonttitle=\bfseries\color{white},
  attach boxed title to top left={yshift=-2mm, xshift=4mm},
  boxed title style={colback=exoheader, colframe=exoheader}]
\textit{Un comprimé d'ibuprofène porte l'indication \og Ibuprofène 400~mg \fg{}.
\textbf{L'indication de l'étiquette est-elle correcte ?}
On se propose de le vérifier par titrage acido-basique avec indicateur coloré et pH-métrie.}
\end{tcolorbox}

\vspace{6pt}

% ══════════════════════════════════════════════════════════════════════════════
\sectitle{Partie I --- Étude préliminaire}
% ══════════════════════════════════════════════════════════════════════════════

\begin{mdframed}[backgroundcolor=exolight, linecolor=exoheader, linewidth=1.2pt]
L'\textbf{ibuprofène} (\ch{C13H18O2}, $M = 206$~g/mol) est un anti-inflammatoire non stéroïdien
comportant un groupe acide carboxylique \ch{-COOH}.
Il est peu soluble dans l'eau mais très soluble dans l'éthanol.
\end{mdframed}

\vspace{4pt}
\question Donner la formule brute et la masse molaire de l'ibuprofène.
Identifier son groupe fonctionnel acide. \dotfill

\vspace{0.3cm}\noindent\dotfill

\question Justifier l'usage de l'éthanol lors de la dissolution du comprimé. \dotfill

\vspace{0.3cm}\noindent\dotfill

\question Écrire l'équation de la réaction de titrage entre l'ibuprofène (\ch{AH})
et les ions hydroxyde (\ch{HO-}). Quelles caractéristiques doit avoir cette réaction ? \dotfill

\vspace{0.3cm}\noindent\dotfill

% ══════════════════════════════════════════════════════════════════════════════
\sectitle{Partie II --- Protocole expérimental}
% ══════════════════════════════════════════════════════════════════════════════

\subsectitle{1. Préparation de la solution d'ibuprofène}

\begin{enumerate}[leftmargin=1.5cm, label=\textcolor{exoheader}{\textbf{\arabic*.}}]
  \item Réduire en poudre le comprimé dans un mortier.
  \item Dissoudre dans 20~mL d'éthanol (les excipients sont insolubles dans l'éthanol).
  \item Filtrer le mélange.
  \item Introduire le filtrat dans un bécher et diluer dans de l'eau distillée.
\end{enumerate}

\vspace{4pt}
\subsectitle{2. Préparation de la burette}

\begin{enumerate}[leftmargin=1.5cm, label=\textcolor{exoheader}{\textbf{\arabic*.}}]
  \item Remplir la burette avec la solution de \ch{NaOH} ($C_b = 0{,}20$~mol/L) et ajuster à 0.
\end{enumerate}

\noindent\textcolor{exoheader}{\faCheckSquare\quad \textbf{Appel n°1} : Faire vérifier le zéro par le professeur.}

\vspace{4pt}
\subsectitle{3. Titrage}

\begin{enumerate}[leftmargin=1.5cm, label=\textcolor{exoheader}{\textbf{\arabic*.}}]
  \item Prélever $V_1 = 10{,}0$~mL de solution d'ibuprofène à la pipette jaugée ; verser dans un erlenmeyer.
  \item Ajouter 10~mL d'eau distillée, quelques gouttes de BBT (indicateur coloré) et le barreau aimanté.
  \item Placer la sonde pH-métrique dans l'erlenmeyer.
  \item Effectuer le titrage en ajoutant la soude par petites fractions ;
        noter le volume $V_E$ correspondant au changement de couleur (jaune → bleu-vert) et au saut de pH.
\end{enumerate}

\vspace{4pt}
\begin{center}
\begin{tabularx}{0.8\linewidth}{|>{\centering\arraybackslash\columncolor{exolight}}X
                                |>{\centering\arraybackslash}X
                                |>{\centering\arraybackslash\columncolor{exolight}}X
                                |>{\centering\arraybackslash}X|}
\hline
\rowcolor{exoheader}
\color{white}$V$ (mL) & \color{white}pH &
\color{white}$V$ (mL) & \color{white}pH \\
\hline
0,0 & & 13,0 & \\
\hline
2,0 & & 13,5 & \\
\hline
4,0 & & 14,0 & \\
\hline
6,0 & & 14,5 & \\
\hline
8,0 & & 15,0 & \\
\hline
10,0 & & 15,5 & \\
\hline
11,0 & & 16,0 & \\
\hline
12,0 & & 17,0 & \\
\hline
12,5 & & 18,0 & \\
\hline
\end{tabularx}
\end{center}

% ══════════════════════════════════════════════════════════════════════════════
\newpage
\sectitle{Partie III --- Exploitation des résultats}
% ══════════════════════════════════════════════════════════════════════════════

\question Tracer la courbe pH~$= f(V)$ et sa courbe dérivée sur papier millimétré.
Déterminer le volume équivalent $V_E$ par la méthode de votre choix.

\[
  V_E = \dotfill \text{ mL}
\]

\question Calculer la quantité de matière d'ions \ch{HO-} versée à l'équivalence,
puis la quantité de matière d'ibuprofène $n_E$.

\vspace{1.5cm}

\question Déduire la masse $m$ d'ibuprofène dans le comprimé et la comparer à 400~mg.

\vspace{1.5cm}

\[
  m = \dotfill \text{ mg}
\]

\question Calculer l'écart relatif. Conclure sur la conformité de l'étiquette. \dotfill

\vspace{0.3cm}\noindent\dotfill

% ══════════════════════════════════════════════════════════════════════════════
\sectitle{Annexe --- Données}
% ══════════════════════════════════════════════════════════════════════════════

\begin{mdframed}[backgroundcolor=exolight, linecolor=exoheader, linewidth=1.2pt]
\begin{itemize}
  \item Ibuprofène : \ch{C13H18O2}, $M = 206$~g/mol, $pK_a = 4{,}9$
  \item \ch{NaOH} : $C_b = 0{,}20$~mol/L (solution titrante dans la burette)
  \item Volume de la prise d'essai : $V_1 = 10{,}0$~mL
  \item Indicateur coloré : BBT (jaune pH < 6 / bleu pH > 7,6)
  \item Phénolphtaléïne : incolore pH < 8,2 / rose pH > 10
\end{itemize}
\end{mdframed}

\end{document}
